%!TEX root = ../thesis.tex

\chapter{Conclusions and Future Work}
\label{chap:conclusions}

\section{Answers to the research questions}

Table~\ref{tab:ch6-axis-ranking} ranks the tested axes under stated conditions,
and Table~\ref{tab:ch7-conclusion-matrix} answers each research question with
its principal boundary.

\begin{table}[H]
\centering
\small
\begin{tabularx}{\textwidth}{@{}p{0.08\textwidth}XX@{}}
\toprule
RQ & Supported conclusion & Principal boundary \\
\midrule
RQ1 & The circularly polarised Alfv\'en wave converged against an exact
solution on both devices and both precisions at pairwise orders rising $1.706$
to $1.886$ towards two, $1.818$ fitted over the ladder, and on the same ladder
fp32 left the fp64 curve at $N=2048$; Brio--Wu approached the same-implementation
HLL $N=8000$ comparator; the Kelvin--Helmholtz linear growth deficit halved
with each doubling towards the published rate; and all cases passed the
declared completion, physical-state and matched-device checks. & Order
verification and the fp32 floor hold for one smooth case, whose pairwise orders
are still rising, so $1.818$ summarises the ladder rather than an asymptotic
order; the Brio--Wu comparator is not exact, and the three-grid 2D comparisons
did not identify a grid range with a stable rate. \\
RQ2 & Responses depended on case and axis: covered HLL CPU/GPU final arrays
agreed only with device contraction disabled, and restoring the compiler
default moved the field further than halving the precision on Brio--Wu;
optimisation level, instruction set and thread count changed nothing, the last
of these on an Euler path measurably running in parallel. & Build comparisons
report density; HLL device evidence covers Brio--Wu and Orszag--Tang. The
non-zero Brio--Wu responses lie inside the p24 Monte Carlo spread, so they are
not separable in order. Timing is one-workstation, five-sample evidence with
different warm-up policies. \\
RQ3 & All 12 same-grid density mean-$L_1$ discrepancies were below 3\% of the
matched fp64 finest adjacent-grid difference; followed in time on the clamped
HLL path at $128^2$ all three cases climb to two or three roundings and stop,
but that plateau belongs to the configuration and not to the arithmetic:
refining doubles the fitted exponent, HLLD quadruples it and removes the
plateau, and the two compound, the series past the roll-up then favouring an
exponential over a power law by a narrow margin. & This
ratio is one refinement-scale diagnostic, not fp32 error or adequacy. The
amplifying window lies past the case's validation stopping time, so it rests on
completion and physical-state checks rather than a reference; loss of
divergence control is excluded as its cause but no positive mechanism is
attributed. Only Brio--Wu has matched deterministic/MCA scope. \\
RQ4 & Configurations, hashes, execution metadata, acceptance criteria and retained summaries
form the reproducibility contract. & Independent-code, cross-machine and MPI
reproduction remain untested. \\
\bottomrule
\end{tabularx}
\caption[Research-question conclusions]{\textbf{Research-question conclusions
and claim boundaries.} Evidence is not transferred between metrics.}
\label{tab:ch7-conclusion-matrix}
\end{table}

\section{Contribution and claim boundaries}

The principal methodological contribution is an experimental design organised
by test case and controlled axis, which keeps numerical-reference error,
cross-variant discrepancy, stochastic spread and runtime separate rather than
treating them as interchangeable measures of ``accuracy''.

Read by case rather than by axis, Orszag--Tang was the most precision-sensitive,
its domain-mean discrepancy reaching 335 roundings at $512^2$ under HLLD;
Brio--Wu was where build and device semantics mattered most, one device flag
there moving the field at least as far as the precision change did;
Liska--Wendroff Configuration~3 was the only case whose fp32 build-semantics
response exceeded its own precision response; and Sod moved for no axis except
precision.

\section{Prioritised future work}

Two priorities follow directly from the limits above. Order verification rests
on one smooth case, so the Alfv\'en wave should be joined by an inclined
two-dimensional problem and by a quadruple-precision Brio--Wu reference
\citep{maddock_boost_multiprecision}, which would turn the present same-scheme
$N=8000$ comparator into an arithmetic ground truth and separate discretisation
error from precision discrepancy without new physics. Second, the
Kelvin--Helmholtz validation should extend past $t=1.0$ so that the window in
which the discrepancy amplifies is covered by a reference rather than by
completion checks alone, with a finite-time Lyapunov exponent of the fp64 flow
measured over the same window: divergence control is excluded as the cause of
that growth (Section~\ref{sec:ch5-saturation-axes}) and nothing yet replaces it.

The remaining mechanisms follow. An HLL variant without the $\pm c_h$ clamp,
solving the GLM pair separately, would separate the Riemann solver from the
wave-speed treatment that makes the tested HLL a globally scaled Rusanov flux.
The $512^2$ positivity failure should be repeated in fp32, and the pressure
recovery replaced by a formulation that does not subtract the dominant energies.
The build matrix should extend beyond the single MSVC environment, after which a
source-level precision-exploration tool \citep{hoeroldEtAl2025raptor} could
search the mixed-precision space that a two-level axis samples only at its
endpoints. Execution coverage should then be made genuinely matched: GPU HLLD
and Kelvin--Helmholtz paths, matched MCA experiments, real OpenMP work sharing
and Message Passing Interface (MPI) reductions with controlled orders, and
repetition on several recorded systems, with kernel, launch, transfer and I/O
costs separated.
